% !PRINTED 521
\weberpage{521}

From this theorem we here draw an important consequence: The functions

\weq{1}{\varphi = \varphi_0 t^h + \varphi_1 t^{h-1} + \cdots + \varphi_h,}

in which the coefficients $\varphi_0, \varphi_1, \ldots, \varphi_h$ are integral or fractional functionals in $\Omega$, but are assumed to be free of the variables $t$, themselves belong to the functionals in the field $\Omega$. Of them the following theorem holds:

2. \gesperrt{A functional $\varphi$ is integral only when the coefficients $\varphi_0, \varphi_1, \ldots, \varphi_h$ are integral functionals.}

To prove it, we assume that $\varphi_0, \varphi_1, \ldots, \varphi_h$ are not all integral at once. If we determine their common denominator $\mu$ and set

\[
\mu\varphi_0 = \alpha_0, \quad \mu\varphi_1 = \alpha_1, \quad \ldots, \quad \mu\varphi_h = \alpha_h,
\]

then $\alpha_0, \alpha_1, \ldots, \alpha_h$ are integral functionals, and $\mu$ contains at least one prime factor $\pi$ which does not at the same time divide all the numerators $\alpha_0,, \alpha_1, \ldots, \alpha_h$ % [sic] doubled comma in print
(cf. the close of the preceding section).

Then the function

\weq{2}{\chi = \mu\varphi = \alpha_0 t^h + \alpha_1 t^{h-1} + \cdots + \alpha_h,}

whose coefficients are now integral, is certainly an integral functional.

If we now assume that $\varphi$ itself is integral, then $\mu\varphi$ is divisible by $\pi$, while yet not all the coefficients $\alpha_0, \alpha_1, \ldots, \alpha_h$ are divisible by $\pi$. Now if $\varphi$ is an integral functional, it satisfies an equation of the form

\weq{3}{\varphi^m = C_1 \varphi^{m-1} + C_2 \varphi^{m-2} + \cdots + C_m,}

in which the coefficients $C_1, C_2, \ldots, C_m$ are integral rational functionals.

Following \S.\,136, (2), we put these functionals in the form

\[
C_1 = \frac{a_1 E_1}{E}, \quad C_2 = \frac{a_2 E_2}{E}, \quad \ldots, \quad C_m = \frac{a_m E_m}{E},
\]

in which $a_1, a_2, \ldots, a_m$ are the absolute values of $C_1, C_2, \ldots, C_m$, hence natural integers (or zero), and $E, E_1, \ldots, E_m$ are primitive integral functions, hence units.

Then equation (3) yields:

\[
E\varphi^m = a_1 E_1 \varphi^{m-1} + a_2 E_2 \varphi^{m-2} + \cdots + a_m E_m,
\]

and by multiplication with $\mu^m$:

\weq{4}{E\chi^m = \mu\left(a_1 E_1 \chi^{m-1} + a_2 E_2 \mu \chi^{m-2} + \cdots + a_m E_m \mu^{m-1}\right).}
